% ============================================================
%  Números Decimais — Apostila Premium | PratiqueJá
%  Engine: XeLaTeX
% ============================================================
\documentclass[12pt]{article}
\usepackage{../../pratiqueja}

% \hlm — destaque de resultado em modo-math (cor sem bold)
\newcommand{\hlm}[1]{{\color{darkblue}#1}}

\setsubject{Números Decimais}

\begin{document}
\pagenumbering{arabic}
\setstretch{1.2}
\color{bodytext}

\heroheader{Números Decimais}%
           {Operações, Comparação e Conversão}%
           {Matemática · Básico}

\setlength{\columnsep}{18pt}
\setlength{\columnseprule}{0.5pt}
\def\columnseprulecolor{\color{exborder}}

\begin{multicols}{2}

%─────────────────────────────────────────────────────────────
\TW{Def}{badgepurple}{Leitura e Escrita}{Valor posicional dos algarismos decimais}

\regra{A \textbf{vírgula} separa a parte inteira da decimal. Cada casa
à direita vale uma fração de potência de $10$:}

\exemp{%
  1ª casa: \textbf{décimos} $(\div 10)$\\
  2ª casa: \textbf{centésimos} $(\div 100)$\\
  3ª casa: \textbf{milésimos} $(\div 1000)$%
}

\exemp{%
  $3{,}5 = 3 + \dfrac{5}{10}$ → ``três vírgula cinco''\\[3pt]
  $2{,}47 = 2 + \dfrac{4}{10} + \dfrac{7}{100}$\\[3pt]
  $0{,}008$ → ``oito milésimos''%
}

%─────────────────────────────────────────────────────────────
\TW{$<\,>$}{darkblue}{Comparação de Decimais}{Comparar posição a posição}

\regra{Compare os algarismos da \textbf{esquerda para a direita},
posição a posição. Complete com zeros à direita para igualar o número
de casas.}

\SH{Exemplo — ordenar $3{,}4$ \ $3{,}45$ \ $3{,}04$}
\exemp{%
  Reescreva: $3{,}40$ \quad $3{,}45$ \quad $3{,}04$\\
  Parte inteira igual $(3)$; compare décimos: $0 < 4$\\[2pt]
  Crescente: $\hlm{3{,}04 < 3{,}40 < 3{,}45}$%
}

%─────────────────────────────────────────────────────────────
\TW{$+\,-$}{darkblue}{Adição e Subtração}{Alinhar as vírgulas}

\regra{Escreva os números \textbf{alinhando as vírgulas}, complete com
zeros e opere coluna a coluna, da direita para a esquerda.}

\SH{Soma — $4{,}5 + 2{,}37$}
\exemp{%
  Complete $4{,}5 = 4{,}50$ e some por coluna:\\[3pt]
  \centering$\begin{array}{r r @{} l}
      & 4 & {,}50\\
    + & 2 & {,}37\\
    \hline
      & 6 & {,}87
  \end{array}$\par\vspace{2pt}
  \raggedright $4{,}5 + 2{,}37 = \hlm{6{,}87}$%
}

\SH{Subtração — $7{,}2 - 3{,}85$ (com empréstimo)}
\exemp{%
  Complete $7{,}2 = 7{,}20$. Onde o dígito de cima é menor, faça
  \textbf{empréstimo}:\\[3pt]
  \centering$\begin{array}{r r @{} l}
      & 7 & {,}20\\
    - & 3 & {,}85\\
    \hline
      & 3 & {,}35
  \end{array}$\par\vspace{2pt}
  \raggedright {\footnotesize Centésimos: $10-5=5$ (empresta);
  décimos: $11-8=3$ (empresta); unidades: $6-3=3$.}\\[2pt]
  $7{,}2 - 3{,}85 = \hlm{3{,}35}$%
}

%─────────────────────────────────────────────────────────────
\TW{$\times$}{darkblue}{Multiplicação de Decimais}{Ignore a vírgula, depois recoloque}

\regra{Multiplique ignorando a vírgula. No resultado, conte da direita
o \textbf{total de casas decimais dos dois fatores} e ponha a vírgula.}

\exemp{%
  \textbf{$1{,}2 \times 0{,}3$}\\
  $12 \times 3 = 36$ \;{\footnotesize($1+1 = 2$ casas)}\\
  Resultado: $\hlm{0{,}36}$%
}

\SH{Multiplicar por 10, 100, 1000…}
\exemp{%
  A vírgula anda para a \textbf{direita}:\\
  $2{,}54 \times 10 = \hlm{25{,}4}$\\
  $2{,}54 \times 100 = \hlm{254}$%
}

%─────────────────────────────────────────────────────────────
\TW{$\div$}{darkblue}{Divisão de Decimais}{Transformar em divisão inteira}

\regra{Multiplique dividendo e divisor pela mesma potência de $10$ para
\textbf{eliminar as casas decimais do divisor}; depois divida.}

\exemp{%
  \textbf{$3{,}6 \div 0{,}4$}\\
  $\times 10$ em ambos: $36 \div 4 = \hlm{9}$\\[3pt]
  \textbf{$0{,}72 \div 0{,}08$}\\
  $\times 100$ em ambos: $72 \div 8 = \hlm{9}$%
}

\nota{Dividir por $10$, $100$, $1000$ move a vírgula para a
\textbf{esquerda}: $25{,}4 \div 10 = 2{,}54$.}

%─────────────────────────────────────────────────────────────
\TW{$\leftrightarrow$}{badgegreen}{Conversão com Frações}{Decimal $\leftrightarrow$ Fração}

\SH{Decimal → Fração}
\regra{Dígitos decimais no numerador, a potência de $10$ no
denominador; depois simplifique.}
\exemp{%
  $0{,}5 = \dfrac{5}{10} = \hlm{\dfrac{1}{2}}$\\[4pt]
  $0{,}75 = \dfrac{75}{100} = \hlm{\dfrac{3}{4}}$\\[4pt]
  $1{,}25 = \dfrac{125}{100} = \hlm{\dfrac{5}{4}}$%
}

\SH{Fração → Decimal}
\regra{Divida o numerador pelo denominador.}
\exemp{%
  $\dfrac{3}{4} = 3 \div 4 = \hlm{0{,}75}$\\[4pt]
  $\dfrac{7}{8} = 7 \div 8 = \hlm{0{,}875}$%
}

\end{multicols}

\end{document}
